This is the placement tutorial: all three seat kinds, internal parts, and ---
deliberately --- two broken designs, because \repl{checkdesign}'s error
messages and the visualizer's red paint teach the model faster than any prose.
By the end you will have built a payload bay: a coupler nested in the airframe
carrying a hollow payload sphere, with a tail section abutted behind.

\begin{note}[Setup]
An empty directory with \code{Aerotech.rse} copied in
(\srcf{scripts/tutorial3.cli} replays everything). Expect this script to exit
with code 1 --- it contains two intentional failures, and the exit code
honestly says so.
\end{note}

\subsection*{A clean three-part stack}

\begin{console}
\muted{# QtRocket CLI - type 'help' for commands, 'quit' to exit.}
\uin{loadmotors Aerotech.rse}
\ok loadmotors: 249 motors from Aerotech.rse
\uin{newdesign NoseCone name=Nose baseRadius=0.02 length=0.12 density=1400}
\ok newdesign: root NoseCone id=2
\uin{addpart Nose BodyTube name=Airframe innerRadius=0.0191 outerRadius=0.02 length=0.40 density=1400}
\ok addpart: BodyTube id=3 under 2
\uin{addpart Airframe FinSet name=Fins finCount=3 rootChord=0.06 tipChord=0.03 span=0.045 sweep=0.03 thickness=0.003 bodyRadius=0.02 density=1400 seat=OnSurface parentStation=0 childStation=0}
\ok addpart: FinSet id=4 under 3
\uin{checkdesign}
\ok checkdesign: 3 parts, contiguous, no overlaps
\end{console}


The same spine as Tutorial~1, at 40\,mm scale: nose, airframe,
fins-on-surface. \repl{checkdesign} passes. Now we break it.

\subsection*{Mistake one: the coupler that pokes through}

A coupler is a short tube nested inside the airframe's bore ---
\code{seat=NestInBore}. For that seat, \code{gap} means \emph{insertion
depth}: how far the child slides in from the parent's fore rim. Here is a
$0.05\m$ coupler inserted only $0.02\m$:

\begin{console}
\uin{addpart Airframe BodyTube name=Coupler innerRadius=0.018 outerRadius=0.019 length=0.05 density=1400 seat=NestInBore gap=0.02}
\ok addpart: BodyTube id=5 under 3
\uin{checkdesign}
\err checkdesign: 1 overlap(s)
  overlap: part 5 (OD 0.019 m) intrudes into part 2 (capacity 0.015 m) by 0.004 m at z=-0.09 m
\uin{listparts}
\ok listparts:
  [2] NoseCone "Nose"  m=0.0703717 kg
    [3] BodyTube "Airframe"  m=0.0619095 kg
      [4] FinSet "Fins"  m=0.025515 kg
      [5] BodyTube "Coupler"  m=0.00813672 kg
  -- composite: mass=0.165933 kg, cg_z=ERR listparts: Part::computeCompositeAt: placement solve failed -- part 5 (OD 0.019 m) intrudes into part 2 (capacity 0.015 m) by 0.004 m at z=-0.09 m
\uin{savedesign payload_broken.qrd}
\ok savedesign: payload_broken.qrd
\end{console}


Three lessons in one failure:

\begin{itemize}
\item The overlap message is \emph{located and quantified}: the $0.03\m$ of
      coupler left sticking out of the bore reaches under the nose skirt at
      $z=-0.09\m$, where its $0.019\m$ outer diameter exceeds the nose's
      $0.015\m$ capacity --- the reported $0.004\m$ intrusion is that
      \emph{radial} penetration, at that axial location. Nothing vague.
\item \repl{listparts} still prints the tree --- geometry exists --- but the
      composite line degrades to an \err{} mid-print: composite mass and CG
      are refused while the placement is unphysical. A broken design cannot
      silently produce plausible-looking numbers.
\item \repl{savedesign} still works. A \code{.qrd} records intent, including
      broken intent --- which lets you save your work mid-repair, and lets
      the visualizer show you the problem
      (\cref{fig:payload-broken}).
\end{itemize}

\screenshot{ug-payload-broken.png}{\code{payload\_broken.qrd} in the
visualizer (Hi-Vis scheme): the offending coupler is painted the fixed error
red, jutting from the airframe's fore rim into the nose skirt. The overlap
diagnostics and this red paint come from the same geometry resolver.}{fig:payload-broken}{0.8}

\subsection*{The fix: insertion depth equals length}

\begin{console}
\uin{removepart 5}
\ok removepart: removed id=5 (BodyTube)
\uin{addpart Airframe BodyTube name=Coupler innerRadius=0.018 outerRadius=0.019 length=0.05 density=1400 seat=NestInBore gap=0.05}
\ok addpart: BodyTube id=6 under 3
\uin{checkdesign}
\ok checkdesign: 4 parts, contiguous, no overlaps
\end{console}


\code{gap=0.05} --- the full coupler length --- seats it flush in the bore.
Note the fix removed the broken part \emph{by id} (\code{removepart 5}); in a
saved script you would prefer names (\cref{sec:parts}).

\subsection*{An internal payload}

\begin{console}
\uin{addpart Coupler HollowSphere name=Payload innerRadius=0.016 outerRadius=0.0175 density=1200 seat=NestInBore gap=0.035}
\ok addpart: HollowSphere id=7 under 6
\uin{checkdesign}
\ok checkdesign: 5 parts, contiguous, no overlaps
\end{console}


A hollow sphere, $17.5$\,mm outer radius, nested inside the coupler. Internal
parts stack the same way external ones do: \seat{NestInBore} into the
coupler's bore, insertion depth placing it fully inside. \repl{checkdesign}
still passes --- the resolver knows a sphere inside a tube inside a tube is
fine, because the bore capacities say so.

\subsection*{Mistake two: the floating tail}

\seat{Abut} means \emph{touching}: child's fore plane against parent's aft
plane, \code{gap} a signed standoff along $+z$. A sign slip ---
\code{gap=-0.02} --- leaves two centimeters of daylight:

\begin{console}
\uin{addpart Airframe BodyTube name=Tail innerRadius=0.0191 outerRadius=0.02 length=0.06 density=1400 seat=Abut gap=-0.02}
\ok addpart: BodyTube id=8 under 3
\uin{checkdesign}
\err checkdesign: 1 air gap(s)
  gap: nothing spans z in [-0.54, -0.52] (0.02 m)
\end{console}


An \emph{air gap} is the other way a stack can be wrong: nothing overlaps, but
the rocket is not one connected body. Same locate-and-quantify style, same
honest refusal to compute a composite.

\subsection*{The fix, and the finished bay}

\begin{console}
\uin{removepart 8}
\ok removepart: removed id=8 (BodyTube)
\uin{addpart Airframe BodyTube name=Tail innerRadius=0.0191 outerRadius=0.02 length=0.06 density=1400 seat=Abut}
\ok addpart: BodyTube id=9 under 3
\uin{checkdesign}
\ok checkdesign: 6 parts, contiguous, no overlaps
\uin{listparts}
\ok listparts:
  [2] NoseCone "Nose"  m=0.0703717 kg
    [3] BodyTube "Airframe"  m=0.0619095 kg
      [4] FinSet "Fins"  m=0.025515 kg
      [6] BodyTube "Coupler"  m=0.00813672 kg
        [7] HollowSphere "Payload"  m=0.00635042 kg
      [9] BodyTube "Tail"  m=0.00928642 kg
  -- composite: mass=0.18157 kg, cg_z=-0.251348 m (t=0)
\uin{setmotor I280DM}
\ok setmotor: I280DM
\uin{setdrag 0.75}
\ok setdrag: 0.75
\uin{setarea 0.001257}
\ok setarea: 0.001257 m^2
\uin{savedesign payload_final.qrd}
\ok savedesign: payload_final.qrd
\uin{launch}
\ok launch: 4506 steps, t_final=45.050s
  apogee    = 2121.606 m @ t=14.450s
  max_speed = 590.492 m/s @ t=1.750s
  downrange = 0.000 m
  landing   = t=45.050s, pos=(0.000, 0.000, -0.483)
CSV /tmp/flight1/qtrocket_run.csv
\uin{quit}
\ok bye
\end{console}


Leave \code{gap} at its default of zero: abut means touching. The finished
six-part design flies (that is Tutorial~4's airframe class) and looks ---
from the outside --- identical to a rocket with no payload at all
(\cref{fig:payload-final}).

\screenshot{ug-payload-final.png}{\code{payload\_final.qrd}: the repaired
stack. The coupler and payload sphere are fully internal --- the outside gives
no hint, which is precisely why \repl{checkdesign}, not the eyeball, is the
arbiter of internal geometry.}{fig:payload-final}{0.8}

\begin{warning}[The red paint and the ERR text are the same fact]
\code{qtrocket-visualizer} runs the \emph{same placement resolver} the
simulator uses. When \repl{checkdesign} names an offender, the visualizer
paints that part red --- one geometry model, two views of its verdict. If the
render looks fine but \repl{checkdesign} objects, believe \repl{checkdesign}:
it can see inside.
\end{warning}

\begin{tip}[What you learned]
\seat{Abut}/\seat{NestInBore}/\seat{OnSurface} and what \code{gap} means to
each; overlap vs air-gap diagnostics, located and quantified; broken designs
save but refuse composites; the visualizer's error-red; and that a script with
intentional failures exits 1 --- which Tutorial~5 turns into a feature.
\end{tip}
